\documentclass{beamer}

\usepackage[portuguese]{babel}
\usepackage[utf8]{inputenc}
\usepackage{minted}
\usepackage{graphicx}
\setbeamertemplate{footline}[frame number]

\title{Ansible para Configuração de Infraestrutura}
\author[João Marcelo Uchôa de Alencar]{João Marcelo Uchôa de Alencar}
\institute{Universidade Federal do Ceará - Quixadá}

\begin{document}
   \begin{frame}
      \titlepage
   \end{frame}

   \begin{frame}
      \frametitle{Introdução}
      \begin{itemize}
      \item Contextualização:
      \begin{itemize}
      \item O Terraform automatiza a criação (provisionamento) de máquinas virtuais.
      \item Uma vez implantadas, as VMs devem ser configuradas (instalação de pacotes, arquivos de conf.).
      \end{itemize}
      \item O que é o Ansible?
      \begin{itemize}
      \item Solução de configuração mantida pela Red Hat.
      \item Arquivos de configuração em YAML ou JSON.
      \item \textbf{Agentless}: Poucas dependências (apenas Python) e usa SSH para acessar as VMs.
      \end{itemize}
      \item Diferença de escopo: Embora o Ansible possa provisionar, focaremos no seu uso para configuração pós-provisionamento.
      \end{itemize}
   \end{frame}

   \begin{frame}
      \frametitle{Conceitos Básicos}
      \begin{itemize}
      \item O Ansible possui duas configurações básicas:
      \begin{itemize}
      \item \textbf{Inventário (Inventory)}: Define quais servidores são alvos das ações (Web, Banco de dados, etc.).
      \item \textbf{Playbook}: Define as tarefas de configuração a serem executadas.
      \end{itemize}
      \item No inventário, utilizamos IP ou nome (DNS) dos servidores.
      \item O Playbook associa tarefas a um servidor específico ou a um grupo.
      \end{itemize}
   \end{frame}

   \begin{frame}[fragile]
      \frametitle{Instalação no Ubuntu}
      \begin{itemize}
      \item Comandos para instalação :
      \end{itemize}
      \footnotesize
      \begin{minted}{bash}
sudo apt-get update
sudo apt-get install software-properties-common
sudo apt-add-repository --yes --update ppa:ansible/ansible
sudo apt install ansible python3-boto3
      \end{minted}
      \normalsize
      \begin{itemize}
      \item \textit{Nota: O pacote python3-boto3 é necessário para interações com a AWS.}
      \end{itemize}
   \end{frame}

   \begin{frame}
      \frametitle{Configuração}
      \begin{itemize}
      \item Arquivo global: \texttt{/etc/ansible/ansible.cfg}.
      \item Boas Práticas:
      \begin{itemize}
      \item Evitar configurações globais para facilitar o compartilhamento de código.
      \item Utilizar variáveis de ambiente.
      \item Manter configurações locais no diretório do projeto (inventários e playbooks).
      \end{itemize}
      \end{itemize}
   \end{frame}

   \begin{frame}
      \frametitle{Tipos de Inventário}
      \begin{itemize}
      \item \textbf{Estático}:
      \begin{itemize}
      \item Servidores listados manualmente em arquivo INI ou YAML.
      \item Uso de IP ou FQDN.
      \end{itemize}
      \item \textbf{Dinâmico}:
      \begin{itemize}
      \item Gerado por \textit{script} ou \textit{plugin}.
      \item Captura informações da nuvem (ex: AWS) em tempo real.
      \item Ideal para ambientes criados via Terraform onde IPs mudam.
      \item Requer AWS CLI configurado.
      \end{itemize}
      \end{itemize}
   \end{frame}

   \begin{frame}[fragile]
      \frametitle{Inventário Estático: Exemplo}
      \begin{itemize}
      \item Exemplo simples (INI) agrupando por função :
      \end{itemize}
      \begin{minted}{ini}
[all:vars]
ansible_ssh_common_args='-o StrictHostKeyChecking=no'
[web]
192.168.0.100
192.168.0.101

[bancodedados]
192.168.0.200
      \end{minted}
      \begin{itemize}
      \item Assume-se que o acesso SSH (chaves) já está configurado no SO para estes IPs.
      \end{itemize}
   \end{frame}

   \begin{frame}[fragile]
      \frametitle{Inventário Estático: Parâmetros}
      \footnotesize
      \begin{minted}{ini}
[web]
servidorweb1 ansible_host=3.81.235.92 ansible_user=ubuntu \
             ansible_private_key_file=~/labsuser.pem
servidorweb2 ansible_host=3.81.235.93 ansible_user=ubuntu \
             ansible_private_key_file=~/labsuser.pem

[bancodedados]
bd1 ansible_host=54.3.97.17 ansible_user=ubuntu \
    ansible_private_key_file=~/labsuser.pem
      \end{minted}
      \normalsize
      \begin{itemize}
      \item Principais parâmetros :
      \begin{itemize}
      \item \texttt{ansible\_host}: IP real do servidor (permite usar alias no nome).
      \item \texttt{ansible\_user}: Usuário SSH.
      \item \texttt{ansible\_private\_key\_file}: Caminho da chave PEM.
      \end{itemize}
      \end{itemize}
   \end{frame}

   \begin{frame}[fragile]
      \frametitle{Verificando Conectividade (Ad-Hoc)}
      \begin{itemize}
      \item O módulo \texttt{ping} verifica o acesso aos servidores listados no inventário .
      \end{itemize}
      \begin{minted}{bash}
# Verificar todos os servidores
ansible -i inventario.ini all -m ping

# Verificar apenas grupo web
ansible -i inventario.ini web -m ping

# Verificar servidor específico
ansible -i inventario.ini servidorweb1 -m ping
      \end{minted}
   \end{frame}

   \begin{frame}[fragile]
      \frametitle{Inventário Dinâmico (AWS)}
      \begin{itemize}
      \item Arquivo YAML (\texttt{aws\_ec2.yml}) que usa plugin para ler tags da AWS .
      \end{itemize}
      \footnotesize
      \begin{minted}{yaml}
plugin: amazon.aws.aws_ec2
aws_profile: default
regions:
  - us-east-1
keyed_groups:
  - key: tags
    prefix: tag
      \end{minted}
      \normalsize
      \begin{itemize}
      \item Comando para visualizar o grafo de recursos:
      \begin{minted}{bash}
ansible-inventory -i aws_ec2.yml --graph
      \end{minted}
      \item Comando para testar a conectividade:
      \begin{minted}{bash}
ansible -i aws_ec2.yml all -m ping \
--ssh-extra-args='-o StrictHostKeyChecking=no' \
-u ubuntu --private-key ˜/labsuser.pem
      \end{minted}
      \end{itemize}
   \end{frame}

   \begin{frame}[fragile]
        \frametitle{Diretório group\_vars}
        \begin{itemize}
        \item Podemos definir variáveis comuns para todos os \textit{hosts}.
        \item Criar diretório \texttt{group\_vars} no mesmo nível.
        \item Colocar as opções em um arquivo chamado all.yml:
        \footnotesize
        \begin{minted}{yaml}
---
# Usuário padrão para conexão
ansible_user: ubuntu

# Caminho para sua chave privada
ansible_ssh_private_key_file: ~/labsuser.pem

# Desativa a verificação de Host Key (evita o prompt de yes/no)
ansible_ssh_extra_args: '-o StrictHostKeyChecking=no'
        \end{minted}
        \end{itemize}
   \end{frame}


   \begin{frame}
      \frametitle{Playbook}
      \begin{itemize}
      \item Escrito em YAML.
      \item Lista tarefas de configuração para grupos de \textit{hosts}.
      \item Vantagem: Uso de \textbf{módulos} abstratos em vez de comandos \textit{bash} puros.
      \item Exemplos de ações (Tasks) :
      \begin{itemize}
      \item Instalar pacotes (\texttt{apt}, \texttt{yum}).
      \item Gerenciar serviços (\texttt{service}, \texttt{systemd}).
      \item Manipular arquivos (\texttt{copy}, \texttt{template}).
      \end{itemize}
      \end{itemize}
   \end{frame}

   \begin{frame}[fragile]
      \frametitle{Exemplo de Playbook: Nginx}
      \begin{minted}{yaml}
---
- hosts: web
  tasks:
    - name: Instalar nginx
      become: true
      apt:
        name: nginx
        state: latest
        update_cache: true

    - name: Iniciar o nginx
      service:
        name: nginx
        state: started
      \end{minted}
      \begin{itemize}
      \item \texttt{become: true}: Executa como sudo.
      \item \texttt{state: latest}: Garante a última versão instalada.
      \end{itemize}
   \end{frame}

   \begin{frame}[fragile]
      \frametitle{Exemplo de Playbook: Múltiplos Grupos}
      \footnotesize
      \begin{minted}{yaml}
- hosts: bancodedados
  tasks:
    - name: Instalar o MySQL Server
      become: true
      apt:
        name: mysql-server
        state: latest
        update_cache: true

    - name: Iniciar o MySQL Server
      service:
        name: mysql
        state: started
      \end{minted}
      \normalsize
      \begin{itemize}
      \item Podemos ter múltiplas entradas de \texttt{hosts} no mesmo arquivo.
      \item Para utilizar os \textit{hosts} da AWS no seu playbook.yml, basta referenciar o nome exato que aparece na hierarquia do grafo (sem o @).
      \end{itemize}
   \end{frame}

   \begin{frame}[fragile]
      \frametitle{Organização: Roles e Diretórios}
      \begin{itemize}
      \item Para projetos complexos, organizamos em \textbf{Roles} (Papéis).
      \end{itemize}
      \footnotesize
      \begin{verbatim}
.
|-- aws_ec2.yml (Inventário Dinâmico)
|-- group_vars
|   `-- all.yml (Opções comuns)
|-- playbook.yml (Orquestrador)
`-- roles
    |-- bancodedados
    |   `-- tasks
    |       `-- main.yml
    `-- web
        `-- tasks
            `-- main.yml
      \end{verbatim}
      \normalsize
      \begin{itemize}
      \item O arquivo \texttt{playbook.yml} passa a apenas referenciar as \textit{roles}:
      \end{itemize}
      \footnotesize
      \begin{minted}{yaml}
- hosts: tag_grupo_web
  roles:
    - web
      \end{minted}
   \end{frame}

   \begin{frame}
      \frametitle{Execução}
      \begin{itemize}
      \item Passos para execução :
      \begin{enumerate}
      \item Configurar credenciais AWS CLI.
      \item Criar infraestrutura (Terraform).
      \item Definir inventário (\texttt{aws\_ec2.yml}).
      \item Executar o comando:
      \end{enumerate}
      \end{itemize}
      \begin{center}
      \texttt{\$ ansible-playbook -i aws\_ec2.yml playbook.yml}
      \end{center}
   \end{frame}

   \begin{frame}
      \frametitle{Saída da Execução}
      \begin{itemize}
      \item \textbf{Gathering Facts}: Ansible verifica acessibilidade e coleta dados do SO.
      \item Estados das Tarefas :
      \begin{itemize}
      \item \textcolor{green}{ok}: Sucesso, nenhuma mudança necessária (idempotência).
      \item \textcolor{yellow}{changed}: Sucesso, alteração realizada no sistema.
      \item \textcolor{red}{unreachable}: Falha na conexão (SSH/Rede).
      \item \textcolor{red}{failed}: Conectou, mas erro ao executar comando (ex: pacote inexistente).
      \end{itemize}
      \item \textbf{Play Recap}: Resumo final de quantos hosts foram alterados ou falharam.
      \end{itemize}
   \end{frame}

   \begin{frame}
      \frametitle{Conclusão}
      \begin{itemize}
      \item Ansible permite descrever infraestrutura de forma declarativa e modular.
      \item A vasta coleção de módulos permite configurar desde pacotes simples até orquestração complexa de clusters.
      \item Integração fluida com Terraform (Terraform cria, Ansible configura).
      \item Documentação oficial: \url{https://docs.ansible.com/}
      \end{itemize}
   \end{frame}

\end{document}